% Self-contained LaTeX edition of vortex_damping_analytical_note.pdf.
% Reconstructed from the analytical note supplied in this conversation;
% typography is not intended to reproduce the original PDF pixel for pixel.
% Compile twice with: pdflatex vortex_damping_analytical_note.tex
% No external figures, bibliography files, or custom fonts are required.
\documentclass[10pt,a4paper]{article}
\usepackage[T1]{fontenc}
\usepackage[utf8]{inputenc}
\usepackage[margin=21mm,top=23mm,bottom=23mm,headheight=14pt,headsep=7mm,footskip=11mm]{geometry}
\usepackage{amsmath,amssymb}
\usepackage{newtxtext,newtxmath}
\usepackage{microtype}
\usepackage{booktabs,array}
\usepackage{xcolor}
\usepackage{tcolorbox}
\usepackage{titlesec}
\usepackage{fancyhdr}
\usepackage[colorlinks=true,linkcolor=navy,citecolor=navy,urlcolor=navy]{hyperref}
\definecolor{navy}{HTML}{183D50}
\definecolor{muted}{HTML}{526574}
\definecolor{pale}{HTML}{F0F5F8}
\definecolor{rulecolor}{HTML}{C8D9E2}
\hypersetup{pdftitle={Why the breathing vortices can damp},
  pdfsubject={An analytical calculation of drift resonance and trapping}}
\pagestyle{fancy}
\fancyhf{}
\fancyhead[L]{\sffamily\small\color{muted}BREATHING VORTICES\quad /\quad ANALYTICAL NOTE}
\fancyhead[R]{\sffamily\small\color{muted}30 September 2026}
\fancyfoot[L]{\sffamily\small\color{muted}Reduced theory; not a fitted simulation result}
\fancyfoot[R]{\sffamily\small\color{muted}\thepage}
\renewcommand{\headrulewidth}{0.4pt}
\setlength{\parindent}{0pt}
\setlength{\parskip}{5pt}
\setlength{\emergencystretch}{2em}
\titleformat{\section}{\large\sffamily\bfseries\color{navy}}{\thesection}{0.6em}{}
\titlespacing*{\section}{0pt}{10pt}{7pt}
\titleformat{\subsection}{\normalsize\sffamily\bfseries\color{navy}}{}{0pt}{}
\titlespacing*{\subsection}{0pt}{8pt}{4pt}
\tcbset{colback=pale,colframe=rulecolor,boxrule=0.45pt,arc=2pt,
  left=8pt,right=8pt,top=6pt,bottom=6pt,before skip=7pt,after skip=7pt}
\newcommand{\dd}{\mathrm{d}}
\newcommand{\ii}{\mathrm{i}}
\newcommand{\E}{\mathrm{e}}
\newcommand{\N}{\mathcal{N}}
\newcommand{\Op}{\mathcal{L}}
\newcommand{\Ham}{\mathcal{H}}
\newcommand{\K}{\mathcal{K}}
\newcommand{\Zfun}{\mathcal{Z}}
\newcommand{\phibar}{\bar\phi}
\newcommand{\varphibar}{\bar\varphi}
\newcommand{\vd}{\bar v_{da}}
\newcommand{\nz}{\mathrm{nz}}
\newcommand{\lin}{\mathrm{lin}}
\newcommand{\tr}{\mathrm{tr}}
\DeclareMathOperator{\PV}{PV}
\DeclareMathOperator{\erfc}{erfc}
\begin{document}
\fontsize{10.5}{13}\selectfont
{\LARGE\sffamily\bfseries\color{navy}Why the breathing vortices can damp\par}
\vspace{3pt}
{\large\sffamily\color{muted}An analytical calculation of drift resonance and trapping\par}
\begin{tcolorbox}
\textbf{Idea.} Calculate phase mixing through particles whose magnetic drift plus zonal advection matches the vortex speed. Then ask whether closed particle orbits suppress that transfer at finite vortex amplitude.\par
\textbf{Status.} The linear response and the reduced integrals below are derived explicitly. Trapping is a candidate explanation for the accelerating decay, not a demonstrated result for the simulation. No numerical lifetime is claimed.
\end{tcolorbox}

\section{Start from the supplied kinetic equations}
The working draft~\cite{draft} uses the orbit-averaged equations
\begin{align}
 \partial_t h_a+\{\phibar,h_a\}+\vd\partial_\alpha h_a
 &=\frac{e_af_0}{T}\bigl(\partial_t\phibar-v_{*a}^{T}\partial_\alpha\phibar\bigr),\label{eq:kinetic}\\
 (1-\lambda_D^2\nabla_\perp^2)\frac{e\phi}{T}
 &=\frac{1}{2n_0}\int(h_p-h_e)\,\dd^3v.\label{eq:field}
\end{align}
Here $e_p=e$, $e_e=-e$ and $\{A,B\}=A_\psi B_\alpha-A_\alpha B_\psi$.
The overbar is the draft's bounce/transit average at fixed energy $\varepsilon$ and pitch angle $\lambda$.

Freeze a slowly evolving, flute-like zonal state $(\Phi,H_a)$ and consider a small travelling perturbation:
\begin{equation}
 \phi=\Phi(\psi)+\varphi(\psi,l)\E^{\ii k\alpha-\ii\omega t},
 \qquad
 h_a=H_a(\psi,\varepsilon,\lambda)+g_a\E^{\ii k\alpha-\ii\omega t}.
 \label{eq:ansatz}
\end{equation}
Real parts are understood. Set $U=\partial_\psi\Phi$; this is an advection velocity in the draft's flux coordinates, not the flux-tube volume also denoted $U$ elsewhere in the draft. The fluctuation need not be flute-like.

Since $\{\Phi,g_a\}=\ii kUg_a$ and $\{\varphibar,H_a\}=-\ii kH_a'\varphibar$, linearisation gives
\begin{equation}
 -\ii[\omega-kU-k\vd]g_a-\ii kH_a'\varphibar
 =-\ii\frac{e_af_0}{T}(\omega+kv_{*a}^{T})\varphibar.
 \label{eq:linear}
\end{equation}
Thus the particle response is
\begin{equation}
 \boxed{\begin{aligned}
 g_a&=\frac{\N_a(\omega)}{\omega-\Omega_a}\,\varphibar,
 &\Omega_a&=k(U+\vd),\\[-1pt]
 &&\N_a(\omega)&=\frac{e_af_0}{T}(\omega+kv_{*a}^{T})-kH_a'.
 \end{aligned}}\label{eq:response}
\end{equation}
The prime means $\partial_\psi$ at fixed $(\varepsilon,\lambda)$. Equation~\eqref{eq:response} is an analytical linear response, not a prescribed decay law. It retains the zonal distribution as well as the zonal potential. Neglecting the mean drive solely because its measured net injection is small would be unjustified.

{\small\color{muted}Basis:~\cite{draft}, kinetic equations and definitions (supplied text, lines 43--62). The frozen-background approximation and subsequent derivations are introduced in this note.\par}

\newpage
\section{Close the field equation and extract the damping}
Define $\mathcal P_k=1-\lambda_D^2\nabla_{\perp,k}^2$, with $\partial_\alpha$ replaced by $\ii k$.
Substituting~\eqref{eq:response} into~\eqref{eq:field} gives a self-consistent mode equation:
\begin{equation}
 \Op(\omega)\varphi\equiv\mathcal P_k\varphi
 -\frac{T}{2n_0e^2}\sum_a e_a\int\dd^3v\,
 \frac{\N_a(\omega)}{\omega-\Omega_a}\varphibar=0.
 \label{eq:operator}
\end{equation}
The velocity integral is evaluated at the field-line position of the field equation; the potential inside the integral is orbit averaged. This keeps the $l$--$\lambda$ coupling. A purely local Doppler-shift estimate does not replace~\eqref{eq:operator}.

\subsection{The resonant contribution}
For the $\E^{-\ii\omega t}$ convention, start in $\operatorname{Im}\omega>0$ and analytically continue the causal response. On the real-frequency boundary,
\begin{equation}
 \frac{1}{\omega_r-\Omega+\ii0}
 =\PV\frac{1}{\omega_r-\Omega}-\ii\pi\delta(\omega_r-\Omega).
 \label{eq:landau}
\end{equation}
Writing $\Op=\Op_R+\ii\Op_I$ gives
\begin{equation}
 \Op_I(\omega_r)\varphi
 =\frac{\pi T}{2n_0e^2}\sum_a e_a\int\dd^3v\,
 \N_a(\omega_r)\varphibar\,
 \delta\!\bigl(\omega_r-k[U+\vd]\bigr).
 \label{eq:imaginary}
\end{equation}
The delta function selects the drift-resonant particles:
\begin{equation}
 \omega_r/k=U(\psi)+\vd(\varepsilon,\lambda).
 \label{eq:resonance}
\end{equation}
Magnetic-drift resonances also occur in the linear dipole theory of Ref.~\cite{mishchenko}; that calculation does not by itself describe a finite-amplitude zonal condensate.

\subsection{Weak-damping formula}
Suppose the principal-value problem has a simple real-frequency mode $\Op_R(\omega_r)\varphi=0$, and the resonant imaginary part is perturbatively small. Write $\omega=\omega_r-\ii\gamma_\lin$, with positive $\gamma_\lin$ denoting amplitude damping. Let $\varphi^\dagger$ be a left null mode of $\Op_R$, and $\langle f,g\rangle=\int f^*g\,\dd^3r$. Expanding~\eqref{eq:operator} and projecting gives
\begin{equation}
 \boxed{\gamma_\lin=
 \frac{\langle\varphi^\dagger,\Op_I(\omega_r)\varphi\rangle}
 {\langle\varphi^\dagger,[\partial_\omega\Op_R(\omega_r)]\varphi\rangle}.}
 \label{eq:gamma}
\end{equation}
All explicit frequency dependence, including that in $\N_a$, is differentiated. In a scalar dispersion relation this is $\gamma_\lin=\mathcal D_I/\mathcal D_R'$. The denominator must be non-zero. The sign must be evaluated: resonance can damp or drive a mode, and negligible net injection is not a proof of positive $\gamma_\lin$.

For an isolated mode of fixed shape,
\begin{equation}
 |\phi_\nz|\propto\E^{-\gamma_\lin t},\qquad
 E_\nz\propto\E^{-2\gamma_\lin t},\qquad
 \Gamma_E\equiv-\partial_t\ln E_\nz=2\gamma_\lin.
 \label{eq:energy-rate}
\end{equation}
Equation~\eqref{eq:gamma} need not describe a strongly damped transient, a continuum response, or a Floquet mode of an appreciably time-periodic background.

\newpage
\section{An explicit analytical velocity integral}
To evaluate the resonance rather than leave it as a formal pole, take a fixed pitch angle and use the quadratic-in-speed magnetic drift in the reduced form
\begin{equation}
 u=\varepsilon/T,\quad \Omega_a=kU+au,\quad
 p(u)=\frac{2}{\sqrt\pi}u^{1/2}\E^{-u},\quad
 \int_0^\infty p(u)\,\dd u=1.
 \label{eq:maxwellian}
\end{equation}
Here $a=k\vd(u=1,\lambda)$ is a signed frequency. The pitch-angle/orbit measure is a separate weight. For this explicit example, approximate $\N_a/f_0$ by $b_0+b_1u$, as when the retained background distribution uses density and temperature moments. This is an additional closure, not a property inferred from the movie.

The relevant energy integral is
\begin{equation}
 J(\omega)=\int_0^\infty
 \frac{p(u)(b_0+b_1u)}{\omega-kU-au+\ii0}\,\dd u.
 \label{eq:J}
\end{equation}
For $a>0$, put $z=(\omega-kU)/a$ and define
$\K_n(z)=\int_0^\infty p(u)u^n/(z-u)\,\dd u$ with the causal continuation. Then
\begin{equation}
 J=\frac{b_0\K_0(z)+b_1\K_1(z)}{a},\qquad
 \K_1=z\K_0-1.
 \label{eq:J-reduced}
\end{equation}
With $u=s^2$, the remaining integral reduces to the plasma dispersion function:
\begin{align}
 \K_0(z)&=\frac{4}{\sqrt\pi}\int_0^\infty
 \frac{s^2\E^{-s^2}}{z-s^2}\,\dd s
 =-2\bigl[1+\sqrt z\,\Zfun(\sqrt z)\bigr],\label{eq:K0}\\
 \Zfun(\zeta)&=\frac{1}{\sqrt\pi}\int_{-\infty}^{\infty}
 \frac{\E^{-s^2}}{s-\zeta}\,\dd s
 =\ii\sqrt\pi\E^{-\zeta^2}\erfc(-\ii\zeta),\qquad
 \operatorname{Im}\zeta>0.\label{eq:plasma-function}
\end{align}
Choose $\sqrt z$ with positive imaginary part and continue causally. The symbol $\Zfun$ here is a special function, not the draft's entropy $Z$. Higher polynomial moments follow from
$\K_{n+1}=z\K_n-\langle u^n\rangle$.

More directly, the absorptive part for either sign of $a$ is
\begin{equation}
 \operatorname{Im}J(\omega_r)
 =-\frac{2\sqrt\pi}{|a|}u_r^{1/2}\E^{-u_r}(b_0+b_1u_r)\Theta(u_r),
 \qquad u_r=\frac{\omega_r-kU}{a}.
 \label{eq:J-imag}
\end{equation}
This assumes real $b_0,b_1$ at $\omega_r$; the right-hand side is understood to vanish for $u_r<0$. For $a<0$, transform the causal prescription as well as $z$; the $a>0$ branch formula cannot be reused blindly. Equation~\eqref{eq:J-imag} makes the resonant population and its sign explicit. Spatial structure, both species and the pitch-angle integral still enter~\eqref{eq:gamma}.

\subsection{A closed-form phase-mixing check, without field feedback}
For comparison only, a passive perturbation obeying $\partial_tg+\ii(kU+au)g=0$ and initially proportional to the Maxwellian has the normalized moment
\begin{equation}
 \begin{aligned}
 C(t)&=\E^{-\ii kUt}\int_0^\infty p(u)\E^{-\ii aut}\,\dd u
 =\E^{-\ii kUt}(1+\ii at)^{-3/2},\\
 |C(t)|^2&=(1+a^2t^2)^{-3/2}.
 \end{aligned}\label{eq:passive}
\end{equation}
This is an exact reduced phase-mixing calculation, but $C$ is a distribution moment, not a self-consistent vortex solution. Its logarithmic decay rate is $3a^2t/(1+a^2t^2)$: it first rises and then falls as $3/t$. Its asymptotic law therefore does not explain accelerating late-time decay; whether the measured interval is asymptotic also has to be checked.

\newpage
\section{The nonlinear question: can trapping suppress mixing?}
The linear response excludes the vortex's modification of particle trajectories. To estimate that effect analytically, use local physical coordinates $(x,y)$ and constant $B_0$, and define the electric-flow streamfunction $\Psi=\phibar/B_0$. These coordinates and velocities are distinct from the flux-coordinate variables in Sections~1--3.

Let $\xi=y-ct$ be the coordinate moving with a nearly steady vortex. For a fixed orbit class, approximate $\vd$ as radially constant across the local region. The perpendicular trajectories satisfy
\begin{equation}
 \dot x=-\partial_\xi\Ham,\quad
 \dot\xi=\partial_x\Ham,\quad
 \Ham=\Psi_z(x)+\Psi_v(x,\xi)+(\vd-c)x.
 \label{eq:hamiltonian}
\end{equation}
Thus closed contours must be sought in the moving, orbit-dependent Hamiltonian, not simply in instantaneous laboratory-frame $\mathbf E\times\mathbf B$ streamlines. This is a local reduction of the advective terms in Ref.~\cite{draft}.

\subsection{Non-zero local shear: the pendulum calculation}
Take a resonant radius $x_r$ satisfying $U_z(x_r)+\vd=c$, where $U_z=\Psi_z'$. If $S_r=U_z'(x_r)\ne0$ and the vortex's radial variation is negligible over the island, expansion gives
\begin{equation}
 \Ham=\text{constant}+\frac{S_r}{2}X^2-A\cos(k_y\xi)+\cdots,
 \quad X=x-x_r,\quad A=|\phibar_v|/B_0.
 \label{eq:pendulum}
\end{equation}
For $S_r>0$ the centre is at $X=\xi=0$. Linearising gives
$\ddot\xi=-k_y^2S_rA\xi$. For $S_r<0$, use the stable point shifted by half a wavelength. The small-orbit frequency and separatrix half-width are therefore
\begin{equation}
 \boxed{\omega_\tr=|k_y|\sqrt{|S_r|A},\qquad
 \Delta x_{\mathrm{half}}=2\sqrt{\frac{A}{|S_r|}}.}
 \label{eq:trapping}
\end{equation}
The width follows by equating the centre-line kinetic term to the potential barrier $2A$. Both quantities decrease as $A^{1/2}$ in this limit.

Trapping can modify a resonant response before linear damping removes the perturbation when $\omega_\tr$ is at least comparable to $\gamma_\lin$. This O'Neil-type comparison is motivated by nonlinear Landau-damping theory~\cite{oneil}, but its application here is a hypothesis. It suggests the crossover estimate
\begin{equation}
 \boxed{A_{\mathrm{crit}}\sim\frac{\gamma_\lin^2}{k_y^2|S_r|}.}
 \label{eq:threshold}
\end{equation}
It is not a derivation of a full amplitude-dependent law $\gamma(A)$. In particular, it does not specify the trapped distribution or its replenishment and erosion by the sinks.

\subsection{At a jet maximum the preceding estimate is singular}
The supplied draft locates the vortices at maxima of the zonal velocity~\cite{draft}. There $U_z'=0$, so inserting the reported maximum shear into~\eqref{eq:threshold} is not valid. Resonances may lie on the jet flanks; otherwise the flow curvature and radial vortex shape must be retained. At an elliptic fixed point of the full $\Ham$, linearised Hamilton's equations give instead
\begin{equation}
 \boxed{\omega_\tr^2=\det(\nabla\nabla\Ham)
 =\Ham_{xx}\Ham_{\xi\xi}-\Ham_{x\xi}^{\,2}>0.}
 \label{eq:hessian}
\end{equation}
This formula includes the vortex's own radial curvature. Its amplitude scaling need not be $A^{1/2}$; an island's existence and its frequency must both be established in the relevant moving frame.

\newpage
\section{Phase mixing is not the irreversible energy sink}
In the draft, magnetic drift and nonlinear advection conserve total free energy $W$. The budget and the nonlinear zonal--vortex exchange are~\cite{draft}
\begin{equation}
 \frac{\dd W}{\dd t}=P-D,\qquad
 (\partial_tE_z)_{\mathrm{NL}}=-(\partial_tE_\nz)_{\mathrm{NL}}.
 \label{eq:budget}
\end{equation}
A decrease in the coherent field can be accompanied by increasingly fine structure in the distribution without immediate destruction of $W$. The simulation's fourth-order hyperdiffusion, rather than physical collisions, supplies irreversible removal. The breathing exchange in~\eqref{eq:budget} cancels when the two field components are added; it is not itself a dissipation term.

\subsection{An exactly solvable model of mixing plus hyperdiffusion}
A second elementary calculation shows how mixing feeds a sink. Near a velocity-space resonance, linearise the orbit frequency as $\Omega(v)\simeq\Omega_r+\beta v$, with $v$ now a local velocity offset. Remove the constant phase and consider the illustrative equation
\begin{equation}
 \partial_tf+\ii\beta vf=-\nu_4\partial_v^4f,\qquad \nu_4>0.
 \label{eq:sink-model}
\end{equation}
This models a local velocity-space sink, not the complete orbit-averaged form of the code's dissipation. With $\widehat f(\eta,t)=\int\E^{-\ii\eta v}f(v,t)\,\dd v$ ($\eta$ here is a Fourier coordinate, not the gradient ratio), it becomes
\begin{equation}
 \partial_t\widehat f-\beta\partial_\eta\widehat f
 =-\nu_4\eta^4\widehat f.
 \label{eq:fourier-sink}
\end{equation}
Along a Fourier characteristic $\eta(t)=\eta_0-\beta t$, the solution is
\begin{align}
 \widehat f(\eta_0-\beta t,t)
 &=\widehat f(\eta_0,0)\exp\!\left[-\nu_4\int_0^t(\eta_0-\beta s)^4\,\dd s\right],\label{eq:characteristic}\\
 \int_0^t(\eta_0-\beta s)^4\,\dd s
 &=\frac{\eta_0^5-(\eta_0-\beta t)^5}{5\beta},\qquad\beta\ne0.
 \label{eq:sink-integral}
\end{align}
For the characteristic launched at $\eta_0=0$,
\begin{equation}
 \boxed{\widehat f(-\beta t,t)=\widehat f(0,0)
 \exp\!\left(-\frac{\nu_4\beta^4t^5}{5}\right),\qquad
 t_{\mathrm{sink}}\sim(\nu_4\beta^4)^{-1/5}.}
 \label{eq:sink-time}
\end{equation}
Mixing makes the velocity-space wavenumber grow as $|\eta|\sim|\beta|t$, and the fourth-order sink becomes stronger as $\eta^4$. This calculation explicitly separates transfer to fine structure from its removal. It applies away from a point where the frequency gradient $\beta$ vanishes.

\begin{tcolorbox}
\textbf{Do not identify~\eqref{eq:sink-time} with the vortex envelope.} It describes a moving Fourier component of the distribution, not $E_\nz(t)$ or even $\widehat f(0,t)$. Nevertheless, it shows that accelerating removal of fine-scale structure can occur without a trapping transition. Acceleration alone does not identify the mechanism.
\end{tcolorbox}
In the sink-free limit, this reduced mixing is reversible at the fine-grained level. The damping of a self-consistent macroscopic mode can still be collisionless; the dependence of the measured slow leak on $\nu_4$ must be established rather than assumed. For $k=0$, the magnetic-drift resonance channel used above disappears. This supports the distinction between the fading non-zonal vortices and the protected flute-like zonal component, without proving indefinite persistence in a more complete physical model.

\newpage
\section{What can be compared with the long run}
The following are the rates supplied in the user's description~\cite{longrun}, not quantities independently extracted from the movie. The conversion uses~\eqref{eq:energy-rate} and applies to a locally exponential envelope of fixed shape.
\begin{center}
\renewcommand{\arraystretch}{1.2}
\begin{tabular}{@{}lll@{}}
\toprule
Time interval & Reported $\Gamma_E$ & Equivalent $\gamma_{\mathrm{amp}}=\Gamma_E/2$\\
\midrule
$4\,000$--$8\,000$   & $5\times10^{-4}$       & $2.5\times10^{-4}$\\
$9\,000$--$9\,500$   & $1.2\times10^{-3}$     & $6\times10^{-4}$\\
$10\,000$--$10\,300$ & $2.4\times10^{-3}$     & $1.2\times10^{-3}$\\
$10\,600$--$11\,000$ & $(4\text{--}7)\times10^{-3}$ & $(2\text{--}3.5)\times10^{-3}$\\
\bottomrule
\end{tabular}
\end{center}
Rates are in the inverse simulation-time units quoted in the draft. They are benchmarks, not predictions of this note.

\subsection{What the analytical work has actually supplied}
Equation~\eqref{eq:response} is the linearised kinetic response about a zonal state. Equations~\eqref{eq:operator}--\eqref{eq:gamma} specify the self-consistent resonance calculation, including its sign. Equations~\eqref{eq:K0} and~\eqref{eq:J-imag} evaluate a Maxwellian energy integral analytically. Equations~\eqref{eq:trapping} and~\eqref{eq:hessian} calculate trapping frequencies in stated local limits. Equations~\eqref{eq:passive} and~\eqref{eq:sink-time} are explicitly solved checks, not replacements for the full field problem.

A numerical value from~\eqref{eq:gamma} requires a specified zonal potential $\Phi$, the velocity-dependent $H_a'$ (or an explicit analytical closure for it), magnetic-drift/orbit coefficients, and a consistent spatial mode and frequency. The potential and box dimensions alone do not fix these. In particular, many distributions can share the same charge moment while differing at the resonances.

The two amplitude-reduced restarts described in Ref.~\cite{longrun} probe whether finite amplitude matters. For a linear problem on the same fixed background, consistently scaling all non-zonal distributions and their self-consistent field gives exactly the same normalised linear evolution. Different normalised decay after such scaling would establish an amplitude-dependent effect, but would not uniquely identify trapping. Moving-frame closed orbits and their resonant populations are the additional physical connection.

\begin{tcolorbox}
\textbf{Conclusion.} Drift-resonant phase mixing provides a concrete candidate channel through which the vortex field can lose energy while the flute-like zonal flow avoids that channel. Nonlinear trapping could suppress it at larger amplitude. The derived response, resonance integral and trapping criteria make this testable; they do not yet establish the observed $\gamma(A)$ or a parameter-free lifetime.
\end{tcolorbox}

\begingroup
\small
\renewcommand{\refname}{Sources and scope}
\begin{thebibliography}{9}
\setlength{\itemsep}{3pt}
\bibitem{draft}
D.~Kennedy and G.~G.~Plunk,
\emph{Explosive instability, quiet plasma: phase-space cascades, coherent vortices, and transport suppression in a dipole pair plasma}.
User-supplied working draft, \texttt{Pasted text(3).txt}.
Basis: kinetic equations (lines 43--62), invariants (64--77), sinks (87), phase mixing (98--100), vortices and breathing (115--124), and free-energy balance (136--150).
\bibitem{longrun}
User-supplied description of the smaller-box, $\eta=1$ run, $0\leq t\leq11\,000$, in a $52\times44\,\lambda_D$ box, with decay-rate estimates and amplitude-reduced restart tests. These newer observations are distinct from the draft's shorter run and are used here as reported.
\bibitem{mishchenko}
A.~Mishchenko, G.~G.~Plunk and P.~Helander,
``Electrostatic stability of electron--positron plasmas in dipole geometry,''
\emph{J. Plasma Phys.} \textbf{84}, 905840201 (2018).
\href{https://doi.org/10.1017/S0022377818000193}{doi:10.1017/S0022377818000193}.
External context for magnetic-drift resonances, not a damping calculation for this condensate.
\bibitem{oneil}
T.~O'Neil, ``Collisionless Damping of Nonlinear Plasma Oscillations,''
\emph{Phys. Fluids} \textbf{8}, 2255--2262 (1965).
\href{https://doi.org/10.1063/1.1761193}{doi:10.1063/1.1761193}.
External context for trapping-modified damping; its result is not transplanted as a proven law here.
\end{thebibliography}
\endgroup
\end{document}
